\documentclass[12pt]{article}
\usepackage[utf8]{inputenc}
\usepackage[T1]{fontenc}
\usepackage{hyperref}
\title{Tjungundji Language Resources for K-12 Curriculum Development}
\author{Compiled by Claude Sonnet 5 (claude-sonnet-5) for Andr\'as Kornai}
\date{}
\begin{document}
\maketitle
\section{Introduction}
Tjungundji (ISO 639-3 tjj) is an Aboriginal language of western Cape York, Queensland, in the Anguthimri group of Paman languages. Historically it was spoken around the Batavia River, Cullen Point and Janie Creek. Later speakers were in the Northern Peninsula Area (New Mapoon, Injinoo, Cowal Creek). My search found no evidence that it is taught in schools. Resources are scarce. One note on the identification: the Glottocode given in the prompt (yang1305) appears to belong to Yangkaal, a different language, so the identification should be checked.
\section{Primary Educational Resources}
None found.
\section{Linguistic and Reference Materials}
\begin{itemize}
\item Wikipedia article on Tjungundji: \url{https://en.wikipedia.org/wiki/Tjungundji}
\item Wikipedia article on Anguthimri, the related language group: \url{https://en.wikipedia.org/wiki/Anguthimri_language}
\end{itemize}
\section{Religious and Community Resources}
None found.
\section{Digital and Audio Resources}
\begin{itemize}
\item Pama Language Centre Tjungundji page, which hosts a ``living library'' (content not checked): \url{https://www.pamacentre.org.au/tjungundji-2/}
\end{itemize}
\section{Curriculum Development Context}
Very little is documented. The Pama Language Centre is the natural first contact for any curriculum work, working with the communities concerned.
\end{document}